\documentclass[12pt, a4paper]{article}
\usepackage[utf8]{inputenc}
\usepackage{graphicx}
\usepackage{geometry}
\usepackage{amsmath}
\usepackage{amssymb}
\usepackage{booktabs}
\usepackage{float}
\usepackage{hyperref}
\usepackage{listings}
\usepackage{xcolor}
\usepackage{caption}
\usepackage{subcaption}

\geometry{margin=1in}

% Code listing style
\definecolor{codegreen}{rgb}{0,0.6,0}
\definecolor{codegray}{rgb}{0.5,0.5,0.5}
\definecolor{codepurple}{rgb}{0.58,0,0.82}
\definecolor{backcolour}{rgb}{0.95,0.95,0.92}

\lstdefinestyle{mystyle}{
    backgroundcolor=\color{backcolour},   
    commentstyle=\color{codegreen},
    keywordstyle=\color{magenta},
    numberstyle=\tiny\color{codegray},
    stringstyle=\color{codepurple},
    basicstyle=\ttfamily\footnotesize,
    breakatwhitespace=false,         
    breaklines=true,                 
    captionpos=b,                    
    keepspaces=true,                 
    numbers=left,                    
    numbersep=5pt,                  
    showspaces=false,                
    showstringspaces=false,
    showtabs=false,                  
    tabsize=2
}
\lstset{style=mystyle}

\title{\textbf{Technical Architecture and Implementation Details of the Adaptive Confidence-Aware and Trust-Based EEG Framework}}

\author{
    \textbf{Mr. M. Pruthvi Raj$^{1, a)*}$, T. Ranjithkumar$^{2, b)}$, G. Harika$^{3, c)}$, D. PriyaChandana$^{4, d)}$, Ch. Babiji$^{5, e)}$} \\
    \vspace{0.2cm}
    $^{1}$Sr. Asst. Prof, Dept. of CSE (AI\&ML), Lakireddy Bali Reddy College of Engg., Mylavaram, A.P, India \\
    $^{2,3,4,5}$Student, Dept. of CSE (AI\&ML), Lakireddy Bali Reddy College of Engg., Mylavaram, A.P, India \\
    \vspace{0.2cm}
    \small $^{a)}$jntupruthvi@gmail.com, $^{b)}$ranjithkumar3742@gmail.com, $^{c)}$harikagajj12@gmail.com \\
    \small $^{d)}$chandudevireddy67@gmail.com, $^{e)}$babjichode7@gmail.com
}
\date{}

\begin{document}

\maketitle

\begin{abstract}
This exhaustive technical document provides the complete architectural, mathematical, and algorithmic foundations for the \textit{Adaptive Confidence-Aware and Trust-Based EEG Framework}. It details the multi-channel data acquisition pipelines, the Hybrid CNN-BiLSTM-Attention deep learning model, and the epistemic uncertainty quantification utilizing Monte Carlo Dropout and Temperature Scaling. Furthermore, this document provides the exact source code implementations, flowcharts, and evaluation metrics required to reproduce the system's high-fidelity Selective Prediction Gate.
\end{abstract}

\tableofcontents
\newpage

\section{Introduction and Motivation}
Modern brain-computer interfaces and neuroergonomic assessments rely heavily on deep neural networks to extract discriminative features from Electroencephalogram (EEG) signals. However, standard deterministic models lack the ability to estimate their own uncertainty, leading to highly confident but incorrect predictions when presented with noisy or out-of-distribution biosignals. This technical report details a safety-critical framework designed to mitigate this issue. 

By integrating Monte Carlo Dropout (to measure epistemic uncertainty) and Temperature Scaling (to calibrate predictive probabilities), we establish a \textbf{Trust Engine}. The engine evaluates a composite Reliability Score $R$ and conditionally abstains from predictions that fall below a safety threshold $\tau$.

\section{System Architecture and Data Pipeline}
\subsection{Signal Acquisition}
The system processes 14-channel EEG streams captured at a sampling rate of 128 Hz. 
The channels utilized are: AF3, F7, F3, FC5, T7, P7, O1, O2, P8, T8, FC6, F4, F8, and AF4.

\subsection{Preprocessing Pipeline}
\begin{enumerate}
    \item \textbf{Windowing:} The continuous stream is segmented into 2-second windows, yielding tensors of shape $(14 \times 256)$.
    \item \textbf{Standardization:} Each channel is normalized using a pre-calibrated `StandardScaler` to achieve $\mu = 0$ and $\sigma = 1$.
\end{enumerate}

\begin{figure}[H]
    \centering
    \includegraphics[width=0.8\textwidth]{fig1.jpg}
    \caption{System Architecture Flowchart depicting the complete pipeline from EEG acquisition to the Selective Prediction Gate.}
    \label{fig:architecture}
\end{figure}

\section{Mathematical Foundations of the Hybrid Model}

\subsection{1D-Convolutional Neural Network (CNN)}
The CNN layers extract spatial and short-term temporal features. For an input sequence $X$, the convolution operation with filter $W$ and bias $b$ is defined as:
\begin{equation}
    H_{conv}^{(t)} = \text{ReLU} \left( \sum_{k=0}^{K-1} W_k \cdot X^{(t-k)} + b \right)
\end{equation}

\subsection{Bidirectional LSTM (BiLSTM)}
The BiLSTM captures long-range temporal dependencies. The LSTM cell updates are governed by:
\begin{align}
    f_t &= \sigma(W_f \cdot [h_{t-1}, x_t] + b_f) \\
    i_t &= \sigma(W_i \cdot [h_{t-1}, x_t] + b_i) \\
    \tilde{C}_t &= \tanh(W_C \cdot [h_{t-1}, x_t] + b_C) \\
    C_t &= f_t * C_{t-1} + i_t * \tilde{C}_t \\
    o_t &= \sigma(W_o \cdot [h_{t-1}, x_t] + b_o) \\
    h_t &= o_t * \tanh(C_t)
\end{align}
The bidirectional aspect concatenates the forward $\overrightarrow{h}_t$ and backward $\overleftarrow{h}_t$ hidden states.

\subsection{Attention Mechanism}
To focus on critical temporal segments, we apply an attention layer:
\begin{align}
    u_t &= \tanh(W_w h_t + b_w) \\
    \alpha_t &= \frac{\exp(u_t^\top u_w)}{\sum_{t'} \exp(u_{t'}^\top u_w)} \\
    s &= \sum_{t} \alpha_t h_t
\end{align}

\begin{figure}[H]
    \centering
    \includegraphics[width=0.8\textwidth]{fig2.jpg}
    \caption{Internal Schematic of the CNN-BiLSTM-Attention Neural Network.}
    \label{fig:model}
\end{figure}

\section{Trust Engine and Uncertainty Estimation}

\subsection{Temperature Scaling Calibration}
To ensure the softmax probabilities reflect true correctness likelihoods, we scale the logits $z$ by a calibrated parameter $T = 0.142$:
\begin{equation}
    \hat{q} = \frac{\exp(z / T)}{\sum_{j} \exp(z_j / T)}
\end{equation}

\subsection{Monte Carlo Dropout}
We execute $M=20$ stochastic forward passes with dropout active.
\begin{align}
    \mu(x) &= \frac{1}{M} \sum_{m=1}^{M} \hat{q}_m(x) \\
    U(x) &= \frac{1}{M} \sum_{m=1}^{M} \left( \hat{q}_m(x) - \mu(x) \right)^2
\end{align}

\subsection{Reliability Composite Score}
The final reliability score $R$ is formulated as:
\begin{equation}
    R = 0.6 \cdot \mu(x) + 0.4 \cdot (1 - \tilde{U}(x))
\end{equation}
where $\tilde{U}(x)$ is the normalized uncertainty.

\section{Implementation Code Repository}

The following sections provide the exact Python implementation driving the Trust Engine.

\subsection{Core Trust Engine Logic}
\begin{lstlisting}[language=Python, caption=Trust Engine Evaluation Method]
import torch
import numpy as np

class SelectivePredictionEngine:
    def __init__(self, model, temperature=1.0, threshold=0.85, mc_passes=20):
        self.model = model
        self.temperature = temperature
        self.threshold = threshold
        self.mc_passes = mc_passes
        
    def evaluate_sample(self, x_tensor):
        self.model.train() # Enable Dropout
        mc_preds = []
        
        with torch.no_grad():
            for _ in range(self.mc_passes):
                logits = self.model(x_tensor)
                scaled_logits = logits / self.temperature
                probs = torch.softmax(scaled_logits, dim=1)
                mc_preds.append(probs.cpu().numpy())
                
        mc_preds = np.array(mc_preds) # Shape: (20, 1, 2)
        mean_probs = np.mean(mc_preds, axis=0)[0]
        variance = np.var(mc_preds, axis=0)[0]
        
        predicted_class = np.argmax(mean_probs)
        confidence = mean_probs[predicted_class]
        uncertainty = np.mean(variance)
        
        # Reliability Calculation
        reliability = (0.6 * confidence) + (0.4 * (1.0 - np.clip(uncertainty * 10, 0, 1)))
        
        accepted = reliability >= self.threshold
        
        return {
            'class': int(predicted_class),
            'confidence': float(confidence),
            'uncertainty': float(uncertainty),
            'reliability': float(reliability),
            'accepted': bool(accepted)
        }
\end{lstlisting}

\subsection{Streamlit Dashboard Integration}
The interactive dashboard provides real-time visualization of the 14-channel waveform alongside the Trust Engine's output.

\begin{lstlisting}[language=Python, caption=Streamlit UI Rendering Logic]
import streamlit as st
import plotly.graph_objects as go

def render_decision_panel(decision, threshold, truth_class):
    r_val = decision['reliability']
    is_accepted = decision['accepted']
    
    color = "#10B981" if is_accepted else "#F59E0B"
    title = "ACCEPTED" if is_accepted else "PREDICTION ABSTAINED"
    
    st.markdown(f"""
    <div style='background: {color}22; border: 1px solid {color}; padding: 20px; border-radius: 12px; text-align: center;'>
        <h2 style='color: {color}; margin: 0;'>{title}</h2>
        <h3 style='margin: 10px 0;'>Reliability R = {r_val:.2f} (Threshold: {threshold:.2f})</h3>
    </div>
    """, unsafe_allow_html=True)
\end{lstlisting}

\begin{figure}[H]
    \centering
    \includegraphics[width=0.8\textwidth]{fig3.jpg}
    \caption{Visual output of the Risk-Coverage empirical evaluation plotting selective accuracy.}
    \label{fig:results}
\end{figure}

\section{Conclusion}
The integration of Monte Carlo Dropout and Temperature Scaling transforms traditional EEG analysis into a highly reliable, mathematically sound framework. The extensive codebase and architectural design detailed in this document prove the viability of deploying such models in real-world neuroergonomic environments.

\end{document}
